\section{Activation patching does not recover the source document}
\label{app:actpatch}

LatentQA~\citep{pan2024latentqa} could not be reproduced from the released code and checkpoint.
Intended as a trained activation decoder baseline, it yields no reported metrics.
The claim we could have drawn is that a trained activation decoder also fails, and a broken baseline produces exactly the appearance our result wants.
We instead report a training-free readout of Qwen3-4B-Instruct hidden states via activation patching~\citep{ghandeharioun2024patchscopes,chen2024selfie}.

For each document, we mean-pool a hidden state $h$ from a chosen layer, insert $h$ into a fixed slot of the completion prompt, and generate tokens.
We sweep layers, slot counts, and decoding style.
The question prompt
\begin{quote}\itshape
Document: ?\\[0.5em]
In one sentence, what is the specific topic of this document?
\end{quote}
reports that the document is empty.
The continuation prompt
\begin{quote}\itshape
The following is a scientific abstract.
\end{quote}
carries topic information.
The sweep selects layer $15$ with eight slots.
This baseline is \texttt{ActPatch}.
Both \doctolora and \texttt{ActPatch} inject a vector into the same model and prompt, differing only in vector origin.
Generated outputs are compared to those from the unpatched prompt to control for prompt prior.

We first test whether the resulting embeddings can identify source documents.
Tested on $500$ abstracts per method retrieving each among $99$ same-subfield distractors, \doctolora retrieved its source $91.8\% \pm 1.2$ (MRR $0.946 \pm 0.008$), while \texttt{ActPatch} managed just $6.0\% \pm 1.1$ (MRR $0.132 \pm 0.011$), which is near chance ($1.0\%$). 
Decodes from activations retain weak document signal. 
Their SBERT similarity to the source is only $0.211$ above baseline ($0.794$ for \doctolora). 
\texttt{ActPatch} decodes also often collapse into repetition ($100$ of $500$, vs $1$ for \doctolora).

We interpolate the same $100$ document pairs as \secref{fusion} by decoding three points along each interpolant ($\alpha=0$, $0.5$, $1$), using a linear mix of full \doctolora embeddings or mean-pooled hidden states for \texttt{ActPatch}.
Each decode $D$ of papers $A$ and $B$ is scored by how much its SBERT similarity (all-mpnet-base-v2, see Appendix~\ref{app:edge-protocol}) to $A$ and $B$ exceeds the baseline similarity between $A$ and $B$.
\begin{equation}
\gamma(D;A,B)=\min\bigl(\mathrm{sim}(D,A),\,\mathrm{sim}(D,B)\bigr)-\mathrm{sim}(A,B).
\label{eq:midpoint-gain}
\end{equation}
A verbatim copy of $A$ or $B$ scores $\gamma=0$.
$\gamma>0$ indicates $D$ is closer to both sources than they are to each other.
$\gamma<0$ means at least one similarity falls below the baseline.

The \doctolora midpoint achieves $\gamma = .171 \pm .011$, positive for 99 of 100 pairs, with endpoints at $-0.020$ and $-0.008$.
The \texttt{ActPatch} midpoint yields $\gamma = -0.039 \pm .015$.
Out of the 50 closest pairs, where the sources are most similar and the required threshold for positivity is highest, \doctolora remains positive for all, while \texttt{ActPatch} does so for 8.

Across 28 PACS node means, \doctolora reaches fuzzy overlap $0.575 \pm 0.043$ while \texttt{ActPatch} gets $0.476 \pm 0.038$. Judge agreement is $0.650 \pm 0.068$ for \doctolora and $0.350 \pm 0.065$ for \texttt{ActPatch}. The \doctolora label is closer on 15 nodes, the activation label wins on 5, and there are 8 ties. 
We note that our results do not eliminate the possibility of hidden state-based decoding strategies.
In fact, we note that the \texttt{ActPatch} is train-free and also small in size compared to \doctolora. 
We instead clarify that \doctolora provides values beyond bare hidden-state-based decoding. 